\begin{proposition}[The maximum as a height inequality]\hypertarget{tag-000Y}{}\label{ari:C.6.4}
\StatementRecord{ari:C.6.4}{proved}{000Y}{\SeriesRef{ari:C.8.5}, \SeriesRef{ari:eq:ch-global}; \citet{demarco2017}; \citet[pp.~183--184]{callSilverman1993}; \citet{silverman2006heights}}

For $t\in K\setminus\{0,1\}$,
\begin{equation}\hypertarget{tag-0038}{}\label{ari:eq:ch-bound}
 \widehat h_O(P_t)\geq
 \frac1{[K:\Q]}\sum_{v\nmid\infty}n_v\lambda_v(P_t)
                    -\frac12\log2.
\end{equation}
The inequality is strict on this smooth base. Over $\Q$ it
is $\widehat h_O(P_t)\geq\log A_t-\frac12\log2$.
The generic height is $\widehat h_O(P)=1/12$, and the
variation divisor is $D_P=(\infty)/12$. Consequently
\[
 \widehat h_O(P_t)=\frac1{12}h(t)+O(1)
 \qquad(t\in\overline\Q\setminus\{0,1\}),
\]
where the constant is uniform over algebraic $t$ and is for
the fixed family and section.
\end{proposition}


\begin{proof}
\begin{enumerate}[label=\arabic*.,leftmargin=*]
\item Apply the analytic maximum theorem \ref{ari:C.8.5} at every embedding
and substitute into \eqref{ari:eq:ch-global}. Equality would
require $j(\sigma(t))=1728$ for every embedding, hence $t=1$,
which has been excluded.

\item For the generic local heights at $0$ and $1$, the minimal
discriminant orders are $8$ and $3$, and the section meets
the nonidentity component of types $IV^*$ and $III$.
Their local corrections are $2/3$ and $1/4$, respectively,
so their normalized local heights are zero. At infinity,
put $s=1/t$, $X=s^{-2}x$, $Y=s^{-3}y$. The new equation is
\[
 y^2=x^3-\tfrac34(1-s)x+\tfrac14(1-s)^2,
 \qquad P=(0,(1-s)/2),\qquad\Delta=27s(1-s)^3.
\]
At $s=0$ the node is $(1/2,0)$, whereas $P=(0,1/2)$ is
nonsingular. Its height contribution is $1/12$.
At all other function-field places the point is integral
and the discriminant is a unit. This proves
$D_P=(\infty)/12$ and the generic height.

\item Theorem~1.1 of \citet[arXiv v2, p.~2]{demarco2017} applies with $K=\Q$,
$B=\mathbb P^1$, the displayed elliptic surface and rational
section, since the generic height is positive. It equips
this divisor with continuous adelic semipositive metrics
whose associated height is exactly $\widehat h_O(P_t)$.
Comparing with the usual height for $(\infty)/12$ gives the
bounded difference. This is the specialization here of the
variation theorem of Tate, with the local refinements of
Silverman reproduced in \citet[Theorems~2.2--2.3, arXiv v2, pp.~7--8]{demarco2017}:
outside finitely many places the local correction is zero,
and at each remaining place it extends continuously after
subtracting its prescribed logarithmic singularity.
Our inequality is an explicit archimedean estimate with
constant $\frac12\log2$.

\item The same generic coefficient appears in
\cite[Theorem~4.1 and the following remark, pp.~183--184]{callSilverman1993}:
for $T=\mathbb P^1$ and $[2]^*(O)\sim4(O)$, the specialization
ratio tends to $1/12$. The bounded remainder used here is the
adelic-height conclusion applied in the preceding step.
\end{enumerate}
\end{proof}